
qui est un
foncteur sur la catégorie \(C\) (grâce à (IV 2.12)). En vertu de (IV 3.2.3), ce dernier commute aux limites inductives filtrantes. D'ailleurs, si \(k'\) est une extension finie de \(k\), alors la projection \(p:X\otimes_k k'\to X\) est un morphisme fini, plat, de présentation finie, donc il définit un homomorphisme image directe \(p=N_{k'/k}:K_\bullet(X\otimes_k k')\longrightarrow K_\bullet(X)\). On vérifie immédiatement que la condition envisagée dans 7.18.3 sur les homomorphismes norme est bien vérifiée. Enfin, si \(A\) est l'algèbre affine d'un ouvert \(U\) de l'espace affine \(\Spec k[T_1,\ldots,T_n]\), alors \(F(k)\to F(A)\) s'identifie à l'application \(K_\bullet(X)\longrightarrow K_\bullet(X\times_k U)\), qui est surjective pour \(X\) noethérien, en vertu de IX 1.6. Utilisant alors les résultats prouvés plus haut, on trouve :
\begin{enumerate}[label=(\roman*)]
\item Si \(K/k\) est une extension pure de \(k\), alors \(K_\bullet(X)\longrightarrow K_\bullet(X_K)\) est injectif, et même bijectif si \(X\) est noethérien.
\item Si \(k\) est séparablement clos et \(K/k\) une extension séparable, alors \(K_\bullet(X)\longrightarrow K_\bullet(X_K)\) est injectif.
\item Pour toute extension \(K\) de \(k\), le noyau de \(K_\bullet(X)\longrightarrow K_\bullet(X_K)\) est un groupe de torsion.
\end{enumerate}
Ces résultats restent valables si on suppose seulement \(X\) noethérien, à condition de se borner à des extensions \(K\) de \(k\) qui sont de type fini (de sorte que \(X_K\) est encore noethérien). De plus, on obtient essentiellement les mêmes résultats pour \(K^\bullet\) au lieu du foncteur \(K_\bullet\), à la seule différence près que dans (i), on ne peut affirmer en général qu'on a bijectivité, faute de disposer du théorème d'homotopie pour \(K^\bullet\) au lieu de \(K_\bullet\); pour pouvoir conclure à la bijectivité, il faudrait supposer \(X\) noethérien régulier, ce qui implique que les groupes \(K^\bullet\) et \(K_\bullet\) sont les mêmes pour \(X\), resp. pour \(X_K\).

